%%% Sample chapter with PDF/A instructions

\chapter{The PDF/A Format}\label{chap:pdfa}

The electronic form of final theses
must be submitted in PDF/A format, level 1a or 2u. These are
profiles of the PDF format that determine which PDF features are permitted
to be used, so that documents are suitable for long-term archiving and
further automated processing. We will focus on level 2u, which we typeset with \TeX.

The most important requirements of PDF/A-2u include:

\begin{itemize}

\item All fonts must be embedded inside the document. References to
external fonts are not permitted (not even to \enquote{system} fonts such as Helvetica or Times).

\item Fonts must contain a ToUnicode table, which defines the mapping from the
character encoding used inside the font to Unicode. This makes it
possible to reliably extract text from the document.

\item The document must contain metadata in XMP format, and if it is in
colour, also a formal specification of the colour space.

\end{itemize}

This template uses the {\tt pdfx} package, which can configure \LaTeX{}
to satisfy the PDF/A requirements. The XMP metadata is generated automatically from the
information in the file matching your main file's name, e.g. {\tt thesis-en.xmpdata} (you can look at the
generated file in {\tt pdfa.xmpi}).

PDF/A validity can be checked using an online validator: \url{https://www.pdf-online.com/osa/validate.aspx/}.

If the file is not valid, common causes include the use of less common
fonts (which are embedded only in bitmap form and/or without Unicode
tables) and the inclusion of images in PDF that do not themselves conform to
the PDF/A standard.

This is likely to affect images created by many different programs.
In such a case, you can try to convert the image to PDF/A using
GhostScript, for example like this:

\begin{verbatim}
        gs -q -dNOPAUSE -dBATCH
           -sDEVICE=pdfwrite -dPDFSETTINGS=/prepress
           -sOutputFile=output.pdf input.pdf
\end{verbatim}
